\documentclass[12pt]{scrartcl}

\usepackage[T1]{fontenc}
\usepackage[utf8]{inputenc}
\usepackage{amsmath}
\usepackage{amsfonts}
\usepackage{amssymb}
\usepackage{amsbsy}
\usepackage{graphicx}
\usepackage{float}
\usepackage{array}
\usepackage{booktabs}
\usepackage{multirow}
\usepackage{bm}
\usepackage{verbatim}
\usepackage[english]{babel}
\usepackage{color}
\usepackage{url}
\usepackage{fancybox}
\usepackage[stable]{footmisc}
\usepackage[format=plain,labelfont=bf]{caption}
\usepackage{fancyhdr}
\usepackage{natbib}
\usepackage{calc}
\usepackage{textcomp}
\usepackage[pdftex,pdfborder={0 0 0}]{hyperref}
\usepackage{pdfpages}
\usepackage{tikz}
\usetikzlibrary{shapes,backgrounds,fit,positioning,arrows,positioning,decorations.pathmorphing,decorations.pathreplacing,backgrounds}
\usepackage{accents}

\renewcommand*{\familydefault}{\sfdefault}

% Margins
\addtolength{\textheight}{1.0cm}
\addtolength{\oddsidemargin}{-0.5cm}
\addtolength{\evensidemargin}{-0.5cm}
\addtolength{\textwidth}{1.0cm}
\parindent=0em

% Appropriate font for \mathcal{}
\DeclareSymbolFont{cmmathcal}{OMS}{cmsy}{m}{n}
\DeclareSymbolFontAlphabet{\mathcal}{cmmathcal}

% Set subscript and superscripts positions
\everymath{
\fontdimen13\textfont2=5pt
\fontdimen14\textfont2=5pt
\fontdimen15\textfont2=5pt
\fontdimen16\textfont2=5pt
\fontdimen17\textfont2=5pt
}

% Bibliography style
\setlength{\bibsep}{1pt}

% Part
\renewcommand\partheadstartvskip{\clearpage\null\vfil}
\renewcommand\partheadmidvskip{\par\nobreak\vskip 20pt\thispagestyle{empty}}
\renewcommand\partheadendvskip{\vfil\clearpage}
\renewcommand\raggedpart{\centering}

\begin{document}
\title{Orographic interpolation}
\author{Benjamin Ménétrier}
\date{Last update: \today}

\thispagestyle{empty}

\maketitle

\section{3D variables}

\subsection{Vertical discretization}

The vertical levels is built with:
\begin{itemize}
\item Half-levels top-down from 0 to $n_z$
\item Full-levels top-down from 0 to $n_z-1$
\end{itemize}

The surface pressure $P_s$ is defined at the last half-level ($n_z$). Hybrid coefficients $a_k$ and $b_k$ are defined at half-levels, to compute the pressure at half-levels:
\begin{align}
P^h_k = a_k + b_k P_s
\end{align}
A full-levels, the temperature $T^f_k$ is known, and the pressure $P^f_k$ is given by:
\begin{align}
P^f_k = \frac{P^h_k + P^h_{k+1}}{2}
\end{align}

\subsection{Height calculation}

The hydrostatic equation is:
\begin{align}
\frac{\partial P}{\partial z} = -\rho g
\end{align}
where:
\begin{itemize}
\item $\rho$ is the density
\item $g$ is the gravitational acceleration
\end{itemize}
and the ideal gas equation is:
\begin{align}
P = \rho R^\star T
\end{align}
where $R^\star$ is the specific gas constant for air. Combining them, we get:
\begin{align}
\label{eq:vertical}
\frac{\partial P}{\partial z} = -\frac{Pg}{R^\star T}
\end{align}
which can be discretized into:
\begin{align}
\delta z =  -\frac{R^\star T}{Pg} \delta P
\end{align}
Integrating this equation upward from the surface, we get iteratively the altitude for each half-level, starting with $z^h_{n_z} = z_s$:
\begin{align}
z^h_{k-1} = z^h_k - \frac{R^\star T^f_{k-1}}{P^f_{k-1} g} \left(P^h_{k-1}-P^h_k\right)
\end{align}
The height at full-levels is approximated by:
\begin{align}
z^f_k = \frac{z^h_{k+1} + z^h_k}{2}
\end{align}
which is slightly inconsistent with the pressure at full levels.

\subsection{Constant height coordinate}

The constant height coordinate is defined by:
\begin{itemize}
\item the constant height surface: $\overline{z}_s = \min(z_s)$
\item the constant height full-levels:
\begin{align}
\overline{z}^f_k = \alpha_k \min(z^f_k) + (1-\alpha_k) \max(z^f_k)
\end{align}
where $\alpha_k$ increases linearly from $\alpha_0 = 0$ to $\alpha_{n_z-1} = 1$.
\end{itemize}
The $n_s$ constant half-levels are defined as:
\begin{subequations}
\begin{align}
\overline{z}^h_{n_s} & = \overline{z}_s \\
\overline{z}^h_k & = \frac{\overline{z}^f_k+\overline{z}^f_{k-1}}{2} \text{ for } 0 < k < n_s \\
\overline{z}^h_0 & = 2 \overline{z}^f_0-\overline{z}^h_{1}
\end{align}
\end{subequations}

\section{Surface pressure}

\subsection{Balanced increment}

Equation \eqref{eq:vertical} can be expressed as:
\begin{align}
\label{eq:vertical_log}
\frac{\partial \log P}{\partial z} = -\frac{g}{R^\star T}
\end{align}
which can be integrated from the constant height surface:
\begin{align}
\log P(z) = \log P(\overline{z}_s) - \int_{\overline{z}_s}^{z} \frac{g}{R^\star T(z')} d z'
\end{align}
If the values of $\log P(\overline{z}_s)$ and $T(z')$ are perturbed with small perturbations $\delta \log P(\overline{z}_s)$ and $\delta T(z')$, then the resulting perturbation on $\log P(z)$ is:
\begin{align}
\delta \log P(z) & = \left(\log P(\overline{z}_s) + \delta \log P(\overline{z}_s)\right) - \int_{\overline{z}_s}^{z} \frac{g}{R^\star \left(T(z')+\delta T(z')\right)} d z' - \log P(z) \nonumber \\
& = \delta \log P(\overline{z}_s) - \int_{\overline{z}_s}^{z} \frac{g}{R^\star} \left(\frac{1}{\left(T(z')+\delta T(z')\right)} - \frac{1}{T(z')}\right) d z'
\end{align}
Assuming that the temperature perturbation is small enough $\delta T(z') \ll T(z')$, then a first-order approximation yields:
\begin{align}
\frac{1}{\left(T(z')+\delta T(z')\right)} - \frac{1}{T(z')} & = \frac{1}{T(z')} \left(\frac{1}{1 + \frac{\delta T(z')}{T(z')}}-1\right) \nonumber \\
 & \simeq -\frac{1}{T(z')} \frac{\delta T(z')}{T(z')} \nonumber \\
& \simeq -\frac{\delta T(z')}{T^2(z')}
\end{align}
Thus:
\begin{align}
\delta \log P(z) & \simeq \delta \log P(\overline{z}_s) + \int_{\overline{z}_s}^{z} \frac{g}{R^\star T^2(z')} \delta T(z') d z'
\end{align}
The integration be can discretized on constant half-levels:
\begin{align}
\delta \log \overline{P}^h_{k} \simeq \delta \log \overline{P}_s + \sum_{k'=k}^{n_s-1} \frac{g}{R^\star \overline{T}^{f2}_{k'}} \left(\overline{z}^h_{k'} - \overline{z}^h_{k'+1}\right) \delta \overline{T}^f_{k'}
\end{align}


\subsection{Interpolation extrapolation}

Assuming a constant lapse rate, the temperature relationship is given by:
\begin{align}
T_2 = T_1-L \left(z_2-z_1\right)
\end{align}
where:
\begin{itemize}
\item $T_1$ and $T_2$ are the temperatures at altitude $z_1$ and $z_2$
\item $L$ is the constant lapse rate
\end{itemize}

Thus, the temperature at constant height full-levels below orography can be computed from the lowest model level:
\begin{align}
\overline{T}^f_k = T^f_{n_z-1} - L \left(\overline{z}^f_k - z^f_{n_z-1}\right)
\end{align}

%\bibliographystyle{mybib-en}
%\bibliography{orographic_interpolation}

\end{document}
